% ===================== CHAPITRE 2 — SPRINT 0 =========================
\chapter{Sprint 0 : Analyse et Spécification des Besoins}

% --- macro locale : logo d'outil non flottant, légendé et numéroté ---
\newcommand{\logofig}[3]{%  #1 = fichier (dossier logos/), #2 = hauteur, #3 = légende
\par\medskip\noindent
\begin{minipage}{\linewidth}\centering
\includegraphics[height=#2]{logos/#1}\par\smallskip
\captionof{figure}{#3}
\end{minipage}\par\medskip}

\section*{Introduction}
Ce chapitre constitue le \emph{Sprint~0}, phase de cadrage qui précède les sprints
de réalisation. Il identifie les acteurs et spécifie les besoins fonctionnels et
non fonctionnels, formalise le pilotage Scrum et le \emph{backlog} produit,
présente les diagrammes de cas d'utilisation global et de classes, planifie les
sprints, puis décrit l'environnement de développement et l'architecture du
système.

% ----------------------------------------------------------------------
\section{Étude du contexte}

\subsection{Identification des acteurs}
Trois profils interagissent avec l'application, conformément aux rôles gérés par
le module de sécurité (table \texttt{SECURITY\_USERS}). Le
tableau~\ref{tab:acteurs} en précise les responsabilités.

\begin{table}[htbp]
\centering
\caption{Identification des acteurs}
\label{tab:acteurs}
\renewcommand{\arraystretch}{1.25}
\begin{tabularx}{\textwidth}{>{\bfseries}p{4cm}X}
\toprule
\textmd{\textbf{Acteur}} & \textbf{Responsabilités}\\
\midrule
Administrateur (\texttt{ADMIN}) & Gère les agences, les employés et les clients ;
valide les demandes de création de clients et d'employés ; consulte le journal
d'audit ; supervise le volet décisionnel.\\
Directeur commercial (\texttt{DIRECTEUR\_COMMERCIAL}) & Soumet les demandes de
création (clients, employés) ; valide les demandes de modification des données
des employés ; consulte les objectifs et les indicateurs.\\
Utilisateur (\texttt{USER}) & Consulte les données autorisées, effectue son
pointage et soumet des demandes de modification de ses propres informations.\\
\bottomrule
\end{tabularx}
\end{table}

\subsection{Spécification des besoins}

\subsubsection{Besoins fonctionnels}
Le système doit permettre :
\begin{itemize}[noitemsep]
  \item l'authentification sécurisée et la gestion des rôles et des droits
        d'accès (\texttt{ADMIN}, \texttt{DIRECTEUR\_COMMERCIAL}, \texttt{USER}) ;
  \item la gestion (CRUD) des agences, des employés et des clients ;
  \item la consultation des objectifs commerciaux et le suivi de la production ;
  \item le \textbf{pointage (présence) des employés} : saisie de l'heure d'arrivée
        et de départ, statut (présent / retard / absent), en mode automatique
        (bouton) ou manuel (administrateur) ;
  \item la soumission et la validation des demandes (clients, employés,
        modifications) selon le cycle \texttt{EN\_ATTENTE} $\rightarrow$
        \texttt{EN\_COURS} $\rightarrow$ \texttt{EFFECTUE}/\texttt{REFUSE} ;
  \item la notification des demandes en attente et la journalisation des actions ;
  \item la consolidation des données par une chaîne \textbf{ETL} et la restitution
        d'indicateurs via un tableau de bord et Power BI ;
  \item la \textbf{segmentation des agences} par apprentissage non supervisé.
\end{itemize}

\subsubsection{Besoins non fonctionnels}
Le tableau~\ref{tab:bnf} récapitule les besoins non fonctionnels.

\begin{table}[htbp]
\centering
\caption{Besoins non fonctionnels}
\label{tab:bnf}
\renewcommand{\arraystretch}{1.2}
\begin{tabularx}{\textwidth}{>{\bfseries}p{3.4cm}X}
\toprule
\textmd{\textbf{Besoin}} & \textbf{Description}\\
\midrule
Sécurité & Authentification, gestion des rôles et protection des opérations
sensibles par des workflows de validation.\\
Fiabilité & Intégrité des données garantie par le SGBD et les transactions JTA.\\
Performance & Temps de réponse acceptable sur les volumes manipulés.\\
Ergonomie & Interface claire et cohérente, navigation par menu latéral.\\
Maintenabilité & Architecture en couches, code modulaire et faiblement couplé.\\
Portabilité & Application Jakarta EE standard, déployable sur serveur compatible.\\
\bottomrule
\end{tabularx}
\end{table}

% ----------------------------------------------------------------------
\section{Pilotage du projet avec Scrum}
Le pilotage du projet repose sur la méthode \textbf{Scrum} présentée au chapitre
précédent. Le tableau~\ref{tab:roles_scrum} précise la répartition des rôles
adoptée durant le stage.

\begin{table}[htbp]
\centering
\caption{Répartition des rôles Scrum du projet}
\label{tab:roles_scrum}
\renewcommand{\arraystretch}{1.2}
\begin{tabularx}{\textwidth}{>{\bfseries}p{3.6cm}p{3.8cm}X}
\toprule
\textmd{\textbf{Rôle}} & \textbf{Assuré par} & \textbf{Mission}\\
\midrule
Product Owner & Maître de stage (BTK) & Exprime et priorise les besoins,
valide les fonctionnalités livrées.\\
Scrum Master & Encadrant académique & Veille au respect de la démarche Scrum et
au bon déroulement des itérations.\\
Équipe de développement & Étudiant & Prend en charge l'analyse, la conception, le
développement et les tests.\\
\bottomrule
\end{tabularx}
\end{table}

% ----------------------------------------------------------------------
\section{Backlog du produit}
Le \emph{product backlog} constitue la liste ordonnée et évolutive des besoins,
exprimés sous forme de \emph{user stories} priorisées selon leur valeur métier.
Il guide la planification des sprints et sert de source unique de travail. Le
tableau~\ref{tab:backlog} en présente le contenu.

\begin{table}[htbp]
\centering
\caption{Backlog produit de l'application}
\label{tab:backlog}
\renewcommand{\arraystretch}{1.2}
\begin{tabularx}{\textwidth}{>{\bfseries}p{0.7cm}p{2.6cm}>{\bfseries}p{1cm}Xp{1.5cm}}
\toprule
\textmd{\textbf{ID}} & \textbf{Fonctionnalité} & \textmd{\textbf{US}} &
\textbf{User story} & \textbf{Priorité}\\
\midrule
01 & Authentification et sécurité & US1 & S'authentifier avec identifiant et mot de passe. & Haute\\
   &                              & US2 & Gérer les rôles et les droits d'accès. & Haute\\
\midrule
02 & Gestion des agences & US3 & Gérer les agences (créer, consulter, modifier, supprimer). & Haute\\
\midrule
03 & Gestion des employés & US4 & Gérer les employés et leur rattachement aux agences. & Haute\\
   &                      & US5 & Gérer les gestionnaires et leurs portefeuilles. & Moyenne\\
\midrule
04 & Gestion des clients & US6 & Gérer les clients (portefeuille). & Moyenne\\
   &                     & US7 & Affecter un client à un gestionnaire. & Moyenne\\
\midrule
05 & Objectifs commerciaux & US8 & Suivre les objectifs commerciaux et la production. & Moyenne\\
\midrule
06 & Pointage & US9 & Pointer la présence (arrivée / départ). & Haute\\
   &          & US10 & Saisir ou corriger un pointage manuel. & Moyenne\\
\midrule
07 & Workflow de demandes & US11 & Soumettre une demande (client, employé, modification). & Moyenne\\
   &                      & US12 & Valider ou refuser une demande. & Moyenne\\
   &                      & US13 & Recevoir des notifications et consulter le journal d'audit. & Moyenne\\
\midrule
08 & Volet décisionnel & US14 & Alimenter le datamart via la chaîne ETL. & Haute\\
   &                   & US15 & Consulter le tableau de bord décisionnel. & Haute\\
   &                   & US16 & Consulter les rapports Power BI. & Moyenne\\
\midrule
09 & Segmentation (IA) & US17 & Segmenter les agences par apprentissage non supervisé. & Moyenne\\
\bottomrule
\end{tabularx}
\end{table}

% ----------------------------------------------------------------------
\section{Diagramme de cas d'utilisation global}
La figure~\ref{fig:uc_global} présente le \textbf{diagramme de cas d'utilisation
global}, qui représente d'un point de vue fonctionnel les interactions entre les
acteurs et les services offerts par le système.

\begin{landscape}
\thispagestyle{plain}
\centering
\vspace*{0.2cm}
\makebox[\linewidth][c]{%
  \includegraphics[width=24cm,height=14.2cm,keepaspectratio]{uc_global.png}}
\par\vspace{0.3cm}
\captionof{figure}{Diagramme de cas d'utilisation global}
\label{fig:uc_global}
\end{landscape}

On y distingue les trois acteurs et leurs cas d'utilisation respectifs. Le cas
\emph{S'authentifier} est commun à tous ; il est relié aux autres par une relation
\texttt{\guillemotleft include\guillemotright}, l'accès aux fonctionnalités exigeant une
authentification préalable. L'\textbf{administrateur} concentre les cas de gestion
et de validation ; le \textbf{directeur commercial} soumet les demandes et consulte
les indicateurs ; l'\textbf{utilisateur} se limite à la consultation, au pointage
et aux demandes le concernant.

% ----------------------------------------------------------------------
\section{Diagramme de classe}
La figure~\ref{fig:class} présente le \textbf{diagramme de classes} du modèle
métier, qui représente les entités persistantes du domaine, leurs attributs et
les associations qui les relient.

\begin{figure}[htbp]
  \centering
  \includegraphics[width=\textwidth]{class.png}
  \caption{Diagramme de classes du modèle métier}
  \label{fig:class}
\end{figure}

L'entité \textbf{\texttt{Utilisateur}} occupe une position centrale. Une
\texttt{Agence} regroupe plusieurs \texttt{Utilisateur}, \texttt{Client} et
\texttt{Objectif} ; un \texttt{Gestionnaire} encadre plusieurs \texttt{Client} ;
les entités de demande (\texttt{DemandeClient}, \texttt{DemandeEmploye},
\texttt{DemandeModification}, \texttt{DemandeAgence}), ainsi que
\texttt{JournalAction}, \texttt{Notification} et \texttt{Pointage}, référencent un
\texttt{Utilisateur} par une association \texttt{@ManyToOne}. Enfin,
\texttt{SecurityUser} est relié à \texttt{Utilisateur} par une dépendance
d'authentification. Le modèle ne comporte ni héritage ni composition : les liens
sont exclusivement des associations, reflétant fidèlement le code du projet.

% ----------------------------------------------------------------------
\section{Planification des sprints}
Conformément au backlog, le développement a été découpé en \textbf{six sprints}
de réalisation, précédés du Sprint~0 de cadrage. Le projet s'est déroulé sur
quatre mois (16 février -- 16 juin 2026). Le tableau~\ref{tab:sprints} détaille
cette planification et la figure~\ref{fig:gantt} en donne le diagramme de Gantt.

\begin{table}[htbp]
\centering
\caption{Planification des sprints}
\label{tab:sprints}
\renewcommand{\arraystretch}{1.2}
\begin{tabularx}{\textwidth}{>{\bfseries}p{1.7cm}p{2cm}p{3.2cm}X}
\toprule
\textmd{\textbf{Sprint}} & \textbf{Durée} & \textbf{User stories} &
\textbf{Objectif du sprint}\\
\midrule
Sprint 1 & 2 semaines & US1, US2 & Authentification et gestion des rôles.\\
Sprint 2 & 2 semaines & US3, US4, US5 & Gestion des agences et des employés.\\
Sprint 3 & 2 semaines & US6, US7, US8 & Gestion des clients et des objectifs.\\
Sprint 4 & 2 semaines & US9--US13 & Pointage, demandes et notifications.\\
Sprint 5 & 3 semaines & US14, US15, US16 & Chaîne ETL et volet décisionnel.\\
Sprint 6 & 2 semaines & US17 & Segmentation des agences (clustering).\\
\bottomrule
\end{tabularx}
\end{table}

\begin{figure}[htbp]
  \centering
  \includegraphics[width=\textwidth]{gantt.png}
  \caption{Planification du projet (diagramme de Gantt des sprints)}
  \label{fig:gantt}
\end{figure}

% ----------------------------------------------------------------------
\section{Environnement de développement}

\subsection{Environnement matériel}
Le développement a été réalisé sur un poste de travail standard : ordinateur
portable doté d'un processeur Intel Core~i5, de 16~Go de mémoire vive et du
système d'exploitation Windows~11.

\subsection{Outil de rédaction du rapport}
Le présent rapport a été rédigé avec \textbf{\LaTeX}, système de composition de
documents particulièrement adapté aux écrits scientifiques et techniques.
\logofig{latex.png}{2.2cm}{Logo de « \LaTeX »}

\subsection{Environnement de développement logiciel}
Plusieurs outils et technologies ont été mobilisés pour la réalisation du projet.

\textbf{Java (JDK~21)} : langage et plateforme d'exécution sur lesquels repose
l'ensemble de l'application.
\logofig{java.png}{2.4cm}{Logo de « Java (JDK 21) »}

\textbf{Jakarta EE~10} : ensemble de spécifications d'entreprise (Servlets, JSP,
EJB, JPA) structurant l'application en couches.
\logofig{jakartaee.png}{2.4cm}{Logo de « Jakarta EE »}

\textbf{Hibernate} : implémentation de JPA assurant le mapping objet-relationnel
entre les entités Java et la base Oracle.
\logofig{hibernate.png}{2.2cm}{Logo de « Hibernate »}

\textbf{Oracle Database} : système de gestion de base de données relationnelle
hébergeant l'ensemble des données opérationnelles et les vues décisionnelles.
\logofig{oracle.png}{1.6cm}{Logo d'« Oracle Database »}

\textbf{WildFly / JBoss} : serveur d'applications Jakarta EE sur lequel
l'application est déployée sous forme d'archive WAR.
\logofig{wildfly.png}{2.2cm}{Logo de « WildFly / JBoss »}

\textbf{Apache Maven} : outil de gestion des dépendances et d'automatisation du
\emph{build} (production du WAR).
\logofig{maven.png}{2cm}{Logo d'« Apache Maven »}

\textbf{IntelliJ IDEA} : environnement de développement intégré utilisé pour le
développement, le débogage et le déploiement de l'application.
\logofig{intellij.png}{2.4cm}{Logo d'« IntelliJ IDEA »}

\textbf{Git et GitHub} : système de gestion de versions et plateforme
d'hébergement du code source.
\par\medskip\noindent
\begin{minipage}{\linewidth}\centering

\includegraphics[height=2.2cm]{logos/git.png}\hspace{1.5cm}%

\includegraphics[height=2.2cm]{logos/github.png}\par\smallskip
\captionof{figure}{Logos de « Git » et « GitHub »}
\end{minipage}\par\medskip

\textbf{ApexCharts} : bibliothèque JavaScript de visualisation utilisée pour les
graphiques du tableau de bord embarqué.
\logofig{apexcharts.png}{2cm}{Logo d'« ApexCharts »}

\textbf{Microsoft Power BI} : outil de \emph{reporting} décisionnel connecté aux
vues analytiques de la base Oracle.
\logofig{powerbi.png}{2.2cm}{Logo de « Microsoft Power BI »}

\textbf{Python} : langage utilisé pour la chaîne \textbf{ETL} et la
\textbf{segmentation} des agences, au moyen des bibliothèques \textbf{pandas}
(manipulation des données) et \textbf{scikit-learn} (apprentissage automatique).
\par\medskip\noindent
\begin{minipage}{\linewidth}\centering

\includegraphics[height=2.2cm]{logos/python.png}\hspace{1.2cm}%

\includegraphics[height=2.2cm]{logos/pandas.png}\hspace{1.2cm}%

\includegraphics[height=2.2cm]{logos/scikitlearn.png}\par\smallskip
\captionof{figure}{Logos de « Python », « pandas » et « scikit-learn »}
\end{minipage}\par\medskip

% ----------------------------------------------------------------------
\section{Architecture du système}
L'architecture définit l'organisation des composants du système et la manière dont
ils communiquent. Elle est présentée à trois niveaux : globale, physique et
logique.

\subsection{Architecture globale}
L'architecture globale (figure~\ref{fig:archi_globale}) offre une vue d'ensemble
des composants. Un poste client (navigateur) dialogue en \texttt{HTTP/HTTPS} avec
l'application web déployée sur \textbf{WildFly}, structurée en couches
(présentation, contrôle, métier, persistance) ; celle-ci accède à la base
\textbf{Oracle} en JDBC. La base alimente, d'une part, \textbf{Power BI} via les
vues \texttt{V\_BI\_*} et, d'autre part, la \textbf{chaîne décisionnelle Python}
(ETL puis segmentation).

\begin{figure}[htbp]
  \centering
  \includegraphics[width=0.92\textwidth]{archi_globale.png}
  \caption{Architecture globale du système}
  \label{fig:archi_globale}
\end{figure}

\subsection{Architecture physique}
L'architecture physique (figure~\ref{fig:archi_physique}) représente le
\textbf{déploiement} des composants sur les nœuds matériels selon une organisation
\textbf{3-tiers} : le poste client, le serveur d'applications WildFly hébergeant
\texttt{gestion-agences.war} (accès à la base via la source de données JTA
\texttt{OracleDS}), le serveur de base de données Oracle (\texttt{FREEPDB1}) et le
poste décisionnel exécutant Power BI.

\begin{figure}[htbp]
  \centering
  \includegraphics[width=0.85\textwidth]{archi_physique.png}
  \caption{Architecture physique (diagramme de déploiement 3-tiers)}
  \label{fig:archi_physique}
\end{figure}

\subsection{Architecture logique}
L'architecture logique (figure~\ref{fig:archi_logique}) décompose le système en
\textbf{couches} faiblement couplées suivant le patron \textbf{MVC} : la couche
\emph{présentation} (JSP/JSTL, ApexCharts), la couche \emph{contrôle} (Servlets et
le filtre \texttt{NotificationFilter}), la couche \emph{métier} (services EJB) et
la couche \emph{persistance} (entités JPA et Hibernate dialoguant avec Oracle via
l'\texttt{EntityManager}).

\begin{figure}[htbp]
  \centering
  \includegraphics[width=0.62\textwidth]{archi_logique.png}
  \caption{Architecture logique en couches de l'application}
  \label{fig:archi_logique}
\end{figure}

\section*{Conclusion}
Ce Sprint~0 a posé les fondations du projet : acteurs, besoins, backlog priorisé,
planification des sprints, environnement de développement et architecture du
système. Les chapitres suivants détaillent chaque sprint de réalisation, en
commençant par l'authentification et la gestion des rôles.
